\SourceLeaf{186}
\SeventhArticle{11}
Wenn die Größen $a, a', a''$ usw. zum Teil oder sämtlich irrational sind, ist es nützlich, die rationale Funktion, aus der wir die Zahlen $R, R', R''$ usw. herleiten, in ein Polynom umzuwandeln. Da sich die Grundsätze einer solchen Umformung nicht in algebraischen Lehrbüchern finden, wollen wir sie hier kurz erläutern. Gegeben seien drei Polynome $Z,\zeta,\zeta'$ in der Unbestimmten $z$; gesucht ist ein Polynom, das an die Stelle des Bruchs $\frac Z\zeta$ treten kann, wenn für $z$ eine beliebige Wurzel der Gleichung $\zeta'=0$ genommen wird. Wir setzen jedoch voraus, dass $\zeta$ für keinen dieser Werte von $z$ verschwindet, oder, was dasselbe bedeutet, dass $\zeta$ und $\zeta'$ keinen gemeinsamen nichtkonstanten Teiler haben. Die Exponenten der höchsten Potenzen von $z$ in $\zeta$ und $\zeta'$ bezeichnen wir mit $k,k'$.

Man dividiere $\zeta$ auf die gewöhnliche Weise durch $\zeta'$, bis der Grad des Restes unter $k'$ liegt. Der Rest sei $=\frac{\zeta''}{\lambda}$, sein Grad $=k''$, sodass $\frac1\lambda z^{k''}$ sein höchstes Glied ist; den Quotienten der Division setzen wir $=\frac p\lambda$. Entsprechend ergebe die Division der Funktion $\zeta'$ durch $\zeta''$ den Rest $\frac{\zeta'''}{\lambda'}$ vom Grad $k'''$ und den Quotienten $\frac{p'}{\lambda'}$; dann ergebe die Division der Funktion $\zeta''$ durch $\zeta'''$ den Rest $\frac{\zeta''''}{\lambda''}$ vom Grad $k''''$ und den Quotienten $\frac{p''}{\lambda''}$, und so fort, bis in der Folge der Funktionen $\zeta'',\ \zeta'''.\ \zeta''''$ usw., die jeweils im höchsten Glied den Koeffizienten $1$ haben, die Funktion $\zeta^{(m)}=1$ erreicht wird. Dass dies schließlich eintreten muss, ist leicht einzusehen, da keine der Funktionen $\zeta,\ \zeta'\ \zeta'',\ \zeta'''$ usw. mit der vorhergehenden einen gemeinsamen nichtkonstanten Teiler haben kann. Somit kann eine Division ohne Rest gewiss nicht eintreten, solange der Grad des Divisors größer als $0$ ist. Wir haben daher die Folge von Gleichungen
\[
\begin{aligned}
\zeta''&=\lambda\zeta-p\zeta'\\
\zeta'''&=\lambda'\zeta'-p'\zeta''\\
\zeta''''&=\lambda''\zeta''-p''\zeta'''\\
\zeta'''''&=\lambda'''\zeta'''-p'''\zeta''''\\
&\text{usw. bis}\\
\zeta^{(m)}&=\lambda^{(m-2)}\zeta^{(m-2)}-p^{(m-2)}\zeta^{(m-1)}
\end{aligned}
\]
Dabei sind $\zeta'',\allowbreak \ \zeta''',\allowbreak \ \zeta''''\ldotp\ldotp\ldotp\ldotp\zeta^{(m)}$ Polynome in $z$ von den Graden $k'',\allowbreak \ k''',\allowbreak \ k''''\ldotp\ldotp\ldotp\ldotp k^{(m)}$; die Zahlen $k',\allowbreak \ k'',\allowbreak \ k'''\ldotp\ldotp\ldotp\ldotp k^{(m)}$ nehmen fortwährend ab bis zur letzten, $k^{(m)}=0$; auch $p,\allowbreak p',\allowbreak p'',\allowbreak p'''$ usw. sind Polynome in $z$ von den Graden $k-k',\allowbreak \ k'-k'',\allowbreak \ k''-k''',\allowbreak \ k'''-k''''$ usw. (ausgenommen den Fall $k<k'$, in dem offenbar $p=0$ zu setzen ist).

Nach diesen Vorbereitungen bilden wir eine \textit{zweite} Folge von Polynomen in $z$, nämlich $\eta,\ \eta',\ \eta'',\ \eta'''\ldotp\ldotp\ldotp\ldotp\eta^{(m)}$. Wir setzen $\eta=1,\ \eta'=0$; alle übrigen\SourceLeaf{187} leiten wir jeweils aus den beiden vorhergehenden nach demselben Gesetz her, durch das die Funktionen $\zeta,\zeta',\zeta'',\zeta'''$ usw. miteinander verknüpft sind, nämlich durch die Gleichungen
\[
\begin{aligned}
\eta''&=\lambda\eta-p\eta'\\
\eta'''&=\lambda'\eta'-p'\eta''\\
\eta''''&=\lambda''\eta''-p''\eta'''\\
\eta'''''&=\lambda'''\eta'''-p'''\eta''''\quad\text{usw. bis}\\
\eta^{(m)}&=\lambda^{(m-2)}\eta^{(m-2)}-p^{(m-2)}\eta^{(m-1)}
\end{aligned}
\]
Offenbar hat hier $\eta''=\lambda$ den Grad $0$; $\eta'''=-\lambda p'$ hat den Grad $k'-k''$, und entsprechend haben die folgenden Funktionen $\eta'''',\eta'''''$ usw. jeweils die Grade $k'-k''',\ k'-k''''$ usw., sodass die letzte, $\eta^{(m)}$, den Grad $k'-k^{(m-1)}$ erreicht.

Betrachten wir ferner die \textit{dritte} Folge von Funktionen, $\zeta-\zeta\eta,\ \zeta'-\zeta\eta',\ \zeta''-\zeta\eta'',\ \zeta'''-\zeta\eta'''$ usw. Zwischen je drei aufeinanderfolgenden Gliedern dieser Folge besteht offenbar eine entsprechende Beziehung, nämlich
\[
\begin{aligned}
\zeta''-\zeta\eta''&=\lambda(\zeta-\zeta\eta)-p(\zeta'-\zeta\eta')\\
\zeta'''-\zeta\eta'''&=\lambda'(\zeta'-\zeta\eta')-p'(\zeta''-\zeta\eta'')\\
\zeta''''-\zeta\eta''''&=\lambda''(\zeta''-\zeta\eta'')-p''(\zeta'''-\zeta\eta''')
\end{aligned}
\]
Für die erste dieser Funktionen ergibt sich $=0$, für die zweite $=\zeta'$; daraus folgt leicht, dass jede durch $\zeta'$ teilbar ist.

Daraus folgt nun ohne Mühe, dass anstelle des Bruchs $\frac Z\zeta$ das Polynom $Z\eta^{(m)}$ genommen werden kann, sofern $z$ nur Werte erhält, die Wurzeln der Gleichung $\zeta'=0$ sind. Denn die Differenz $\frac{Z(1-\zeta\eta^{(m)})}{\zeta}$ muss für einen solchen Wert von $z$ verschwinden, da $1-\zeta\eta^{(m)}=\zeta^{(m)}-\zeta\eta^{(m)}$ durch $\zeta'$ teilbar ist.

Anstelle der Funktion $Z\eta^{(m)}$ kann auch der bei ihrer Division durch $\zeta'$ entstehende Rest genommen werden, dessen Grad kleiner ist als der Grad der Funktion $\zeta'$.

Dieser Rest lässt sich jedoch durch den folgenden Algorithmus bequemer unmittelbar bestimmen. Man bilde die Gleichungen
\[
\begin{aligned}
Z&=q'\zeta'+Z'\\
Z'&=q''\zeta''+Z''\\
Z''&=q'''\zeta'''+Z'''\\
Z'''&=q''''\zeta''''+Z''''\quad\text{usw. bis}\\
Z^{(m-1)}&=q^{(m)}\zeta^{(m)}+Z^{(m)}
\end{aligned}
\]
\SourceLeaf{188}indem man der Reihe nach $Z$ durch $\zeta'$, dann den Rest der ersten Division, $Z'$, durch $\zeta''$, dann den Rest der zweiten Division durch $\zeta'''$ dividiert und so fortfährt. Da der Rest stets einen niedrigeren Grad hat als der Divisor, sind die Grade der Funktionen $Z',Z'',Z''',Z''''$ usw. jeweils kleiner als $k',k'',k''',k''''$ usw. Die letzte, $Z^{(m)}$, wird notwendig $=0$, da der Divisor $\zeta^{(m)}$ $=1$ ist. Wir haben daher
\[
Z=q'\zeta'+q''\zeta''+q'''\zeta'''+q''''\zeta''''+\text{usw.}+q^{(m)}\zeta^{(m)}
\]
Wenn für $z$ nur Wurzeln der Gleichung $\zeta'=0$ genommen werden, ist $\zeta'=0$, $\zeta''=\zeta\eta''$, $\zeta'''=\zeta\eta'''$, $\zeta''''=\zeta\eta''''$ usw.; unter derselben Einschränkung gilt daher
\[
\frac{Z}{\zeta}=q''\eta''+q'''\eta'''+q''''\eta''''+\text{usw.}+q^{(m)}\eta^{(m)}
\]
Der Grad dieses Ausdrucks muss unter $k'$ liegen. Da nämlich die Grade der Quotienten $q'',q''',q''''$ usw. unter $k'-k''$, $k''-k'''$, $k'''-k''''$ usw. liegen müssen, sind die Grade der einzelnen Teile $q''\eta''$, $q'''\eta'''$, $q''''\eta''''$ usw. kleiner als $k'-k''$, $k'-k'''$, $k'-k''''$ usw.

Schließlich bemerken wir noch: Sollten unter den Werten der Unbestimmten $z$, die in den Bruch $\frac{Z}{\zeta}$ einzusetzen sind, rationale und irrationale Werte gemischt vorkommen, so ist es vorteilhafter, jene von diesen zu trennen und nur diese in der Gleichung $\zeta'=0$ zusammenzufassen. Für rationale Werte ist eine solche Verkürzung der Rechnung nicht nötig; für irrationale Werte wird die Rechnung dagegen umso einfacher, je niedriger der Grad des Polynoms ist, auf das sich die rationale Funktion zurückführen lässt.

\SeventhArticle{12}
Hier folgt nun ein Beispiel der im vorigen Artikel erläuterten Umformung. Gegeben sei die rationale Funktion
\[
\frac{z^6-\frac{50}{39}z^4+\frac{283}{715}zz-\frac{256}{15015}}{7z^6-\frac{105}{13}z^4+\frac{315}{143}zz-\frac{35}{429}}
\]
in der $z$ eine beliebige Wurzel der folgenden Gleichung bezeichnet:
\[
z^7-\frac{21}{13}z^5+\frac{105}{143}z^3-\frac{35}{429}z=0
\]
Sollten hier alle sieben Wurzeln einbezogen werden, so würden wir zu einem Polynom sechsten Grades gelangen. Für den rationalen Wert $z=0$ ist die Berechnung des Bruchs jedoch offenbar unmittelbar möglich und liefert den Wert $\frac{256}{1225}$. Lassen wir diese Wurzel also beiseite, so beschränken wir uns auf die Gleichung sechsten Grades:\SourceLeaf{189}
\[
z^6-\frac{21}{13}z^4+\frac{105}{143}zz-\frac{35}{429}=0
\]
Auf diese Weise können wir leicht voraussehen, dass ein Polynom vierten Grades entsteht. Die Anwendung der vorhergehenden Vorschriften liefert nun Folgendes:
\[
\begin{aligned}
\zeta&=7z^6-\frac{105}{13}z^4+\frac{315}{143}zz-\frac{35}{429}\\[5pt]
\zeta'&=z^6-\frac{21}{13}z^4+\frac{105}{143}zz-\frac{35}{429}\\[5pt]
\zeta''&=z^4-\frac{10}{11}zz+\frac{5}{33}\\[5pt]
\zeta'''&=zz-\frac{3}{7}\\[5pt]
\zeta''''&=1
\end{aligned}\]
\[\begin{aligned}
\lambda&=\frac{13}{42}&p&=\frac{13}{6}\\[5pt]
\lambda'&=-\frac{4719}{280}&p'&=-\frac{4719}{280}zz+\frac{3333}{280}\\[5pt]
\lambda''&=-\frac{147}{8}&p''&=-\frac{147}{8}zz+\frac{777}{88}
\end{aligned}\]
\[\begin{aligned}
\eta&=1\\[5pt]
\eta'&=0\\[5pt]
\eta''&=\frac{13}{42}\\[5pt]
\eta'''&=\frac{20449}{3920}zz-\frac{14443}{3920}\\[5pt]
\eta''''&=\frac{61347}{640}z^4-\frac{127413}{1120}zz+\frac{120263}{4480}
\end{aligned}\]
\[\begin{aligned}
Z&=z^6-\frac{50}{39}z^4+\frac{283}{715}zz-\frac{256}{15015};&q'&=1\\[5pt]
Z'&=\frac{1}{3}z^4-\frac{22}{65}zz+\frac{323}{5005}&q''&=\frac{1}{3}\\[5pt]
Z''&=-\frac{76}{2145}zz+\frac{632}{45045}&q'''&=-\frac{76}{2145}\\[5pt]
Z'''&=-\frac{4}{3465}&q''''&=-\frac{4}{3465}
\end{aligned}
\]
Schließlich ergibt sich das folgende Polynom, das dem gegebenen Bruch äquivalent ist:
\[
-\frac{1859}{16800}z^4-\frac{1573}{29400}zz+\frac{7947}{39200}
\]

